\par\Needspace{6\baselineskip}\medskip
\begin{center}\hypertarget{native-246-article175}{}\pdfbookmark[2]{175.}{article-175}175.\end{center}
Ad hunc finem a problemate nostro speciali ad disquisitionem generalissimam
in omni. calculi ad philosophiam naturalem applicatione foecundissimam
ascendemus. Sint $V$, $V'$, $V''$ etc. functiones incognitarum $p$, $q$, $r$,
$s$ etc., $\mu$ multitudo illarum functionum, $\nu$ multitudo incognitarum,
supponamusque, per observationes immediatas valores functionum ita inventos
esse $V=M$, $V'=M'$, $V''=M''$ etc. Generaliter itaque loquendo evolutio
valorum incognitarum constituet problema indeterminatum, determinatum, vel
plus quam determinatum, prout fuerit $\mu<\nu$, $\mu=\nu$, vel
$\mu>\nu$\footnote{Si in casu tertio functiones $V$, $V'$, $V''$ etc. ita
comparatae essent, ut $\mu+1-\nu$ ex ipsis vel plures tamquam functiones
reliquarum spectare liceret, problema respectu harum functionum etiamnum
plus quam determinatum foret, respectu quantitatum $p$, $q$, $r$, $s$ etc.
autem indeterminatum: harum scilicet valores ne tunc quidem determinare
liceret, quando valores functionum $V$, $V'$, $V''$ etc. absolute exacti
dati essent: sed hunc casum a disquisitione nostra excludemus.}.
Hic de ultimo tantum casu sermo erit, in quo manifesto exacta cunctarum
observationum repraesentatio tunc tantum possibilis foret, ubi illae omnes
ab erroribus absolute immunes essent. Quod quum in rerum natura locum non
habeat, omne systema valorum incognitarum $p$, $q$, $r$, $s$ etc. pro possibili
habendum erit, ex quo valores \SourceLeaf{247}functionum $M-V$, $M'-V'$,
$M''-V''$ etc. oriuntur, limitibus errorum, qui in istis observationibus
committi potuerunt, non maiores: quod tamen neutiquam ita intelligendum est,
ac si singula haec systemata possibilia aequali probabilitatis gradu gauderent.

Supponemus primo, eum rerum statum fuisse in omnibus observationibus, ut
nulla ratio adsit, cur aliam alia minus exactam esse suspicemur, sive ut
errores aeque magnos in singulis pro aeque probabilibus habere oporteat.
Probabilitas itaque cuilibet errori $\Delta$ tribuenda exprimetur per functionem
ipsius $\Delta$, quam per $\varphi\Delta$ denotabimus. Iam etiamsi hanc
functionem praecise assignare non liceat, saltem affirmare possumus, eius
valorem fieri debere maximum pro $\Delta=0$, plerumque aequalem esse pro
valoribus aequalibus oppositis ipsius $\Delta$, denique evanescere, si pro
$\Delta$ accipiatur error maximus vel maior valor. Proprie itaque
$\varphi\Delta$ ad functionum discontinuarum genus referre oportet, et si
quam functionem analyticam istius loco substituere ad usus practicos nobis
permittimus, haec ita comparata esse debebit, ut utrimque a $\Delta=0$
asymptotice quasi ad $0$ convergat, ita ut ultra istum limitem tamquam vere
evanescens considerari possit. Porro probabilitas, errorem iacere inter
limites $\Delta$ et $\Delta+\mathrm{d}\Delta$ differentia infinite parva
$\mathrm{d}\Delta$ ab invicem distantes, exprimenda erit per
$\varphi\Delta\, .\,\mathrm{d}\Delta$; proin generaliter probabilitas,
errorem iacere inter $D$ et $D'$, exhibebitur per integrale
$\int\varphi\Delta\, .\,\mathrm{d}\Delta$, a $\Delta=D$ usque ad $\Delta=D'$
extensum. Hoc integrale a valore maximo negativo ipsius $\Delta$ usque ad
valorem maximum positivum, sive generalius a $\Delta=-\infty$ usque ad
$\Delta=+\infty$ sumtum, necessario fieri debet $=1$.

Supponendo igitur, systema aliquod determinatum valorum quantitatum $p$,
$q$, $r$, $s$ etc. locum habere, probabilitas, pro $V$ ex observatione
proditurum esse valorem $M$, exprimetur per $\varphi(M-V)$, substitutis
in $V$ pro $p$, $q$, $r$, $s$ etc. valoribus suis; perinde
$\varphi(M'-V')$, $\varphi(M''-V'')$ etc. expriment probabilitates, ex
observationibus resultaturos esse functionum $V'$, $V''$ etc. valores $M'$,
$M''$ etc. Quamobrem quandoquidem omnes observationes tamquam eventus ab
invicem independentes spectare licet, productum
\[
\varphi(M-V)\, .\,\varphi(M'-V')\, .\,\varphi(M''-V'')\, .\,\text{etc.}=\Omega
\]
exprimet exspectationem seu probabilitatem, omnes istos valores simul ex
observationibus prodituros esse.
